\documentclass[11pt]{article}
\usepackage[margin=1in]{geometry}
\usepackage{amsmath,amssymb,amsthm}
\usepackage{graphicx}
\usepackage[hidelinks]{hyperref}
\setlength{\emergencystretch}{2em}

\newtheorem{theorem}{Theorem}
\newtheorem{lemma}{Lemma}
\newtheorem{corollary}{Corollary}

\newcommand{\C}{\mathbb C}
\newcommand{\M}{\mathcal M}
\newcommand{\omegaone}{\omega_1}

\title{Non-Sigma-Unital Coronas Are Universal for Unital \(C^*\)-Algebras}
\author{Run \texttt{fa\_banach\_001}, lane 19}
\date{June 17, 2026}

\begin{document}
\maketitle

\section*{Source Problem}

The source paper is Jinghao Huang, Karimbergen Kudaybergenov, and Rui Liu,
\emph{The ball-covering property of von Neumann algebras and noncommutative
symmetric spaces}, arXiv:2311.14257.  After proving that a commutative unital
\(C^*\)-algebra \(C(K)\) has BCP/UBCP exactly when \(K\) has a countable
\(\pi\)-basis, the authors ask for a noncommutative version:

\begin{center}
\includegraphics[width=0.92\textwidth]{figures/open_problem_crop.png}
\end{center}

\section*{Statement}

The previous packet
\[
        \texttt{2311.14257\_sigma\_unital\_coronas\_fail\_bcp}
\]
proved that if \(A\) is nonunital and \(\sigma\)-unital, then
\(\M(A)/A\) fails BCP.  The following theorem shows that the
\(\sigma\)-unital hypothesis is essential.

\begin{theorem}
\label{thm:universal-corona}
Let \(B\) be a nonzero unital \(C^*\)-algebra.  There is a nonunital
non-\(\sigma\)-unital \(C^*\)-algebra \(A\) such that
\[
        \M(A)/A\cong B.
\]
Indeed, one can take
\[
        A=C_0([0,\omega_1),B).
\]
\end{theorem}

\begin{corollary}
\label{cor:positive-corona}
There is a nonunital non-\(\sigma\)-unital \(C^*\)-algebra \(A\) such that
\(\M(A)/A\) has UBCP.  For example,
\[
        A=C_0([0,\omega_1))
\]
satisfies
\[
        \M(A)/A\cong \C.
\]
\end{corollary}

\section*{Idea of the Proof}

The ordinal space \([0,\omega_1)\), with the order topology, is locally compact
and non-\(\sigma\)-compact.  Its decisive feature here is that every continuous
map from \([0,\omega_1)\) to a metric space is eventually constant.  Since a
\(C^*\)-algebra is a metric space in norm, every bounded continuous
\(B\)-valued function on \([0,\omega_1)\) has an eventual value in \(B\).
For unital \(B\),
\[
        \M(C_0([0,\omega_1),B))=C_b([0,\omega_1),B),
\]
and the ideal \(C_0([0,\omega_1),B)\) is precisely the kernel of the eventual
value map.  The corona is therefore \(B\).

\section*{Proof}

Let
\[
        X=[0,\omega_1)
\]
with its order topology.

\begin{lemma}
\label{lem:eventually-constant}
If \(Y\) is a metric space and \(f:X\to Y\) is continuous, then \(f\) is
eventually constant.
\end{lemma}

\begin{proof}
Fix \(\varepsilon>0\).  We first show that there is
\(\alpha_\varepsilon<\omega_1\) such that
\[
        \operatorname{diam} f([\alpha_\varepsilon,\omega_1))<\varepsilon.
\]
If not, start with \(\alpha_0=0\).  Having chosen \(\alpha_n\), choose
ordinals \(\beta_n,\gamma_n\ge\alpha_n\) such that
\[
        d(f(\beta_n),f(\gamma_n))\ge\varepsilon.
\]
Then choose \(\alpha_{n+1}\) strictly larger than both \(\beta_n\) and
\(\gamma_n\).
Let
\[
        \delta=\sup_n \alpha_n.
\]
This supremum is still a countable ordinal, hence \(\delta<\omega_1\).  Since
both \(\beta_n\) and \(\gamma_n\) eventually lie above every \(\alpha_N\), both
sequences converge to \(\delta\) in the order topology.  By continuity,
\(f(\beta_n)\to f(\delta)\) and \(f(\gamma_n)\to f(\delta)\), contradicting
\[
        d(f(\beta_n),f(\gamma_n))\ge\varepsilon
\]
for all \(n\).

Thus for every \(m\ge1\) there is \(\alpha_m<\omega_1\) such that
\[
        \operatorname{diam} f([\alpha_m,\omega_1))<1/m.
\]
Let \(\alpha=\sup_m\alpha_m\).  Again \(\alpha<\omega_1\).  If
\(\xi,\eta\ge\alpha\), then \(d(f(\xi),f(\eta))\le 1/m\) for every \(m\), so
\(f(\xi)=f(\eta)\).  Hence \(f\) is constant on \([\alpha,\omega_1)\).
\end{proof}

\begin{lemma}
\label{lem:corona-identification}
Let \(B\) be a unital \(C^*\)-algebra.  Then
\[
        \M(C_0(X,B))/C_0(X,B)\cong B.
\]
\end{lemma}

\begin{proof}
Since \(B\) is unital, the standard multiplier identification gives
\[
        \M(C_0(X,B))=C_b(X,B),
\]
with pointwise operations and norm.  By Lemma~\ref{lem:eventually-constant},
each \(f\in C_b(X,B)\) is eventually constant.  Define
\[
        \rho:C_b(X,B)\longrightarrow B
\]
by letting \(\rho(f)\) be this eventual value.  The map \(\rho\) is a
surjective \(*\)-homomorphism, since every \(b\in B\) is the eventual value of
the constant function \(x\mapsto b\).

Its kernel is exactly \(C_0(X,B)\).  Indeed, if \(f\in C_b(X,B)\) has eventual
value \(0\), then \(f\) is eventually zero, so its support is contained in
\([0,\alpha]\) for some \(\alpha<\omega_1\); this interval is compact, hence
\(f\in C_0(X,B)\).  Conversely, if \(f\in C_0(X,B)\) had nonzero eventual value
\(b\ne0\), then for \(0<\varepsilon<\|b\|\) the set
\[
        \{x\in X:\|f(x)\|\ge\varepsilon\}
\]
would contain a tail \([\alpha,\omega_1)\), which is not compact.  This
contradicts \(f\in C_0(X,B)\).  Therefore
\[
        C_b(X,B)/C_0(X,B)\cong B.
\]
\end{proof}

\begin{lemma}
\label{lem:not-sigma-unital}
If \(B\ne0\) is unital, then \(A=C_0(X,B)\) is nonunital and not
\(\sigma\)-unital.
\end{lemma}

\begin{proof}
The space \(X=[0,\omega_1)\) is locally compact and noncompact, so
\(C_0(X,B)\) is nonunital.

We next recall that \(X\) is not \(\sigma\)-compact.  Every compact subset of
\([0,\omega_1)\) is bounded by some countable ordinal.  A countable union of
such bounded compact sets is still bounded below \(\omega_1\), and therefore
cannot cover \(X\).

Suppose, toward a contradiction, that \(C_0(X,B)\) were \(\sigma\)-unital.
Then it has a countable positive contractive approximate unit \((u_n)\).  The
compact sets
\[
        K_n=\{x\in X: \|u_n(x)\|\ge 1/2\}
\]
cover \(X\).  Indeed, fix \(x\in X\).  Choose \(g\in C_c(X)\) with
\(g(x)=1\) and choose \(b\in B\) with \(\|b\|=1\).  Put
\[
        a(t)=g(t)b\in C_0(X,B).
\]
For some \(n\), \(\|u_na-a\|<1/2\).  Evaluating at \(x\) gives
\[
        \|u_n(x)b-b\|<1/2,
\]
and hence \(\|u_n(x)\|\ge1/2\).

Since each function \(x\mapsto\|u_n(x)\|\) belongs to \(C_0(X)\), each
\(K_n\) is compact.  Thus \(X=\bigcup_n K_n\) would be \(\sigma\)-compact, a
contradiction.  Therefore \(A\) is not \(\sigma\)-unital.
\end{proof}

\begin{proof}[Proof of Theorem~\ref{thm:universal-corona}]
Take \(A=C_0(X,B)\).  Lemma~\ref{lem:not-sigma-unital} says that \(A\) is
nonunital and not \(\sigma\)-unital.  Lemma~\ref{lem:corona-identification}
gives
\[
        \M(A)/A\cong B.
\]
\end{proof}

\begin{proof}[Proof of Corollary~\ref{cor:positive-corona}]
Apply Theorem~\ref{thm:universal-corona} with \(B=\C\).  The one-dimensional
space \(\C\) has UBCP, so the corona \(\M(C_0([0,\omega_1)))/C_0([0,\omega_1))\)
has UBCP.
\end{proof}

\section*{Consequences for Question 3}

The sigma-unital corona theorem is sharp in a strong sense.  Sigma-unital
nonunital coronas all fail BCP by the annulus construction, but
non-\(\sigma\)-unital coronas do not have a uniform BCP behavior: they can be
made isomorphic to any prescribed unital \(C^*\)-algebra.  Therefore the
remaining arbitrary \(C^*\)-algebra classification cannot be reduced to a
blanket statement that coronas fail BCP.

\section*{Verification and Search Bounds}

The run indexes were searched for arXiv:2311.14257 together with
non-\(\sigma\)-unital corona, multiplier corona, \([0,\omega_1)\), ordinal
space, eventual constancy, BCP, and UBCP.  Existing packets contained the
sigma-unital corona failure theorem and the earlier unital stabilized-corona
obstruction, but no packet recorded this non-\(\sigma\)-unital universality
phenomenon.

The proof uses only standard facts about \(C_0(X,B)\), multiplier algebras for
unital coefficient algebras, and the elementary topology of the ordinal space
\([0,\omega_1)\).

\section*{Human Review Focus}

Check Lemma~\ref{lem:eventually-constant}, especially the diagonal sequence
contradiction for tail diameters, and Lemma~\ref{lem:corona-identification},
where \(C_0(X,B)\) is identified as the kernel of the eventual-value map.
Those two lemmas carry the theorem.

\begin{thebibliography}{9}

\bibitem{HKL}
Jinghao Huang, Karimbergen Kudaybergenov, and Rui Liu,
\emph{The ball-covering property of von Neumann algebras and noncommutative
symmetric spaces},
arXiv:2311.14257, 2023.

\bibitem{LLLZ}
Minzeng Liu, Rui Liu, Jimeng Lu, and Bentuo Zheng,
\emph{Ball covering property from commutative function spaces to
non-commutative spaces of operators},
arXiv:2111.04921, 2021.

\end{thebibliography}

\end{document}
